\kbvsechead{ಕೃದಂತ ಭಾವನಾಮ :}

\indent ಕ್ರಿಯೆಯ ಭಾವವನ್ನು ತಿಳಿಸುವ ಶಬ್ದಗಳಿಗೆ ಕೃದಂತ ಭಾವನಾಮಗಳೆಂದು ಹೆಸರು; ಧಾತುಗಳ ಮೇಲೆ ಭಾವಾರ್ಥದಲ್ಲಿ ಕೃತ್‌ ಪ್ರತ್ಯಯಗಳು ಸೇರಿ ಕೃದಂತ ಭಾವನಾಮಗಳಾಗುವುದು.\par

\indent ಉದಾಹರಣೆಗೆ, “ಕುಣಿ” ಎಂಬ ಧಾತುವಿನ ಮೇಲೆ “ತ” ಎಂಬ ಪ್ರತ್ಯಯ ಸೇರಲು “ಕುಣಿತ” ಎಂದಾಗುವುದು. “ಕುಣಿತ” ಎಂಬ ಶಬ್ದರೂಪವು ಕುಣಿಯುವ ಭಾವವನ್ನು ತಿಳಿಸುವುದು. ಆದುದರಿಂದ ಇದು ಕೃದಂತ ಭಾವನಾಮವೆನಿಸುವುದು.\par

\begin{kbvrule}
\noindent\kbvbullet\hspace{0.5em}ಕೆಳಗೆ ಕೊಟ್ಟಿರುವ ಉದಾಹರಣೆಗಳಿಂದ ಕೃದಂತ ಭಾವನಾಮಗಳ ಅಥವಾ ಭಾವ ಕೃದಂತಗಳ ಅರ್ಥ ಮತ್ತಷ್ಟು ಸ್ಪಷ್ಟವಾಗಿ ತಿಳಿದುಬರುವುದು.
\end{kbvrule}

\vspace{5pt plus 4pt minus 1pt}
\noindent{\renewcommand{\arraystretch}{1.24}\setlength{\tabcolsep}{4pt}\rowcolors{2}{kbvrow}{white}
\begin{tabularx}{\textwidth}{X X X}
{\kbvsecfont\color{kbvaccentdark}ಮೂಲ} & {\kbvsecfont\color{kbvaccentdark}ಭಾವಾರ್ಥದಲ್ಲಿ} & {\kbvsecfont\color{kbvaccentdark}ಭಾವ} \\[1pt]
\arrayrulecolor{kbvline}\specialrule{0.5pt}{0pt}{2pt}\arrayrulecolor{black}
ಧಾತು & ಕೃತ್‌ ಪ್ರತ್ಯಯ & ಕೃದಂತ \\
ನಡುಗು &  + ಉಕ & ನಡುಕ \\
ನಗು &  + ಎ & ನಗೆ \\
ನೋಡು &  + ಟ & ನೋಟ \\
ಸೆಳೆ &  + ತ & ಸೆಳೆತ \\
ತೊಡು &  + ಇಗೆ & ತೊಡಿಗೆ \\
ತೊಡು &  + ಗೆ & ತೊಡುಗೆ \\
ನಾಚು &  + ಅಕೆ & ನಾಚಿಕೆ \\
ಉಳು &  + ವುದು & ಉಳುವುದು \\
\end{tabularx}}
\vspace{5pt plus 4pt minus 1pt}

\kbvsechead{ಕೃದಂತಾವ್ಯಯ :}

\indent ಧಾತುವಿನಿಂದ ಹುಟ್ಟಿ ಅವ್ಯಯದ ಲಕ್ಷಣಗಳನ್ನು ಹೊಂದಿರುವ ಶಬ್ದರೂಪಗಳಿಗೆ ಕೃದಂತಾವ್ಯಯಗಳು ಎಂದು ಹೆಸರು.\par

\indent ಧಾತುಗಳ ಮೇಲೆ ಅಲು, ಅಲಿಕ್ಕೆ, ದು, ಇ, ದರೆ, ಉತ್ತ್ಯ, ಅದೆ, ಅ, ಉತ್ತ್ಯ, ಇರಲು-ಮೊದಲಾದ ಪ್ರತ್ಯಯಗಳು ಸೇರಲು ಕೃದಂತಾವ್ಯಯಗಳು.\par
